\documentclass[12pt, a4paper]{article}

% ── Pacotes ───────────────────────────────────────────────────────────────────
\usepackage[utf8]{inputenc}
\usepackage[T1]{fontenc}
\usepackage[brazil]{babel}
\usepackage{geometry}
\usepackage{graphicx}
\usepackage{booktabs}
\usepackage{tabularx}
\usepackage{multirow}
\usepackage{array}
\usepackage{xcolor}
\usepackage{colortbl}
\usepackage{hyperref}
\usepackage{amsmath}
\usepackage{listings}
\usepackage{enumitem}
\usepackage{float}
\usepackage{caption}
\usepackage{subcaption}
\usepackage{setspace}
\usepackage{fancyhdr}
\usepackage{titlesec}
\usepackage{mdframed}
\usepackage{tcolorbox}
\tcbuselibrary{most}

\geometry{
  top=2.5cm, bottom=2.5cm,
  left=3cm, right=2cm
}

% ── Cores ─────────────────────────────────────────────────────────────────────
\definecolor{ufpiblue}{RGB}{0, 51, 102}
\definecolor{lightblue}{RGB}{230, 240, 255}
\definecolor{lightgreen}{RGB}{220, 240, 220}
\definecolor{lightred}{RGB}{255, 230, 230}
\definecolor{lightyellow}{RGB}{255, 250, 220}
\definecolor{codegray}{RGB}{245, 245, 245}
\definecolor{successgreen}{RGB}{0, 128, 0}
\definecolor{warningred}{RGB}{180, 0, 0}

% ── Estilo de código ──────────────────────────────────────────────────────────
\lstset{
  backgroundcolor=\color{codegray},
  basicstyle=\ttfamily\footnotesize,
  breaklines=true,
  frame=single,
  rulecolor=\color{gray!40},
  xleftmargin=1em,
  xrightmargin=1em,
  aboveskip=0.5em,
  belowskip=0.5em
}

% ── Cabeçalho/rodapé ──────────────────────────────────────────────────────────
\pagestyle{fancy}
\fancyhf{}
\rhead{\textcolor{gray}{\footnotesize Questão 4 — Destilação de Conhecimento}}
\lhead{\textcolor{gray}{\footnotesize UFPI -- Tópicos em IA -- 2026.1}}
\rfoot{\thepage}
\renewcommand{\headrulewidth}{0.4pt}

% ── Título de seções ──────────────────────────────────────────────────────────
\titleformat{\section}{\large\bfseries\color{ufpiblue}}{}{0em}{\thesection\quad}[\titlerule]
\titleformat{\subsection}{\normalsize\bfseries\color{ufpiblue!80}}{}{0em}{\thesubsection\quad}

% ── Caixas coloridas ──────────────────────────────────────────────────────────
\newtcolorbox{destaquebox}[1][]{
  colback=lightblue, colframe=ufpiblue,
  fonttitle=\bfseries\small, title=#1,
  left=6pt, right=6pt, top=4pt, bottom=4pt
}
\newtcolorbox{successbox}[1][]{
  colback=lightgreen, colframe=successgreen!70,
  fonttitle=\bfseries\small, title=#1,
  left=6pt, right=6pt, top=4pt, bottom=4pt
}
\newtcolorbox{exemplobox}[1][]{
  colback=lightyellow, colframe=gray!50,
  fonttitle=\bfseries\footnotesize\color{ufpiblue}, title=#1,
  left=6pt, right=6pt, top=4pt, bottom=4pt,
  fontupper=\footnotesize
}
\newtcolorbox{warningbox}[1][]{
  colback=lightred, colframe=warningred!70,
  fonttitle=\bfseries\small, title=#1,
  left=6pt, right=6pt, top=4pt, bottom=4pt
}

% ═════════════════════════════════════════════════════════════════════════════
\begin{document}
% ═════════════════════════════════════════════════════════════════════════════

% ── Capa ──────────────────────────────────────────────────────────────────────
\begin{titlepage}
  \centering
  \vspace*{1cm}
  {\large Questão 4 - Destilação de Conhecimento em Modelos de Linguagem}\\
  \rule{\linewidth}{0.5pt}\\
  \vspace{1.5cm}
  \begin{tabular}{ll}
    \textbf{Modelos}    & Qwen2.5-3B-Instruct (Professor) $\rightarrow$ Qwen2.5-1.5B-Instruct (Aluno) \\[0.3cm]
    \textbf{Técnica}    & Chain-of-Thought Distillation + LoRA (SFTTrainer) \\[0.3cm]
    \textbf{Dataset}    & DOMPI-2025 (Diários Oficiais dos Municípios do Piauí) \\[0.3cm]
    \textbf{Benchmark}  & 100 perguntas sobre documentos municipais piauienses \\
  \end{tabular}
  \vfill
\end{titlepage}

% ── Sumário ───────────────────────────────────────────────────────────────────
\newpage
\newpage

% ═════════════════════════════════════════════════════════════════════════════
\section{Introdução}
% ═════════════════════════════════════════════════════════════════════════════


A destilação de conhecimento (\textit{Knowledge Distillation}), proposta originalmente
por \textit{Hinton et al.} (2015), é uma técnica que permite transferir o conhecimento
de um modelo maior (professor) para um modelo menor (aluno), tornando o aluno capaz de
aproximar o comportamento do professor com menos parâmetros e menor custo computacional.

\subsection{LLMs Usados para Destilação}

Modelos usados para a questão atual:

\begin{itemize}
  \item \textbf{Qwen2.5-3B $\rightarrow$ Qwen2.5-1.5B} -- escolha deste trabalho
\end{itemize}

A escolha de modelos da \textbf{mesma família} é recomendada pois compartilham o mesmo
vocabulário e tokenizador, facilitando a destilação por logits e reduzindo problemas de
incompatibilidade de distribuições.

% ═════════════════════════════════════════════════════════════════════════════
\section{Configuração Experimental}
% ═════════════════════════════════════════════════════════════════════════════

\subsection{Modelos}

\begin{table}[H]
\centering
\caption{Modelos utilizados no experimento de destilação}
\begin{tabular}{llrrr}
\toprule
\textbf{Papel} & \textbf{Modelo} & \textbf{Parâmetros} & \textbf{VRAM (fp16)} & \textbf{Quant.} \\
\midrule
Professor & Qwen/Qwen2.5-3B-Instruct & 3B & $\sim$6 GB & fp16 \\
Aluno     & Qwen/Qwen2.5-1.5B-Instruct & 1.5B & $\sim$3 GB & fp16 \\
\bottomrule
\end{tabular}
\end{table}

Ambos os modelos pertencem à família \textbf{Qwen2.5} (Alibaba Cloud), compartilhando
a mesma arquitetura Transformer com RoPE, SwiGLU e GQA. A versão \texttt{-Instruct}
foi escolhida por ter passado por ajuste fino supervisionado (SFT) e alinhamento por
RLHF, sendo adequada para tarefas de perguntas e respostas sobre documentos.

\subsection{Dataset Sintético}

O dataset de destilação foi gerado sinteticamente a partir do corpus
\textbf{DOMPI-2025}.

\textbf{Processo de geração:}
\begin{enumerate}
  \item Seleção de 100 trechos representativos do corpus DOMPI-2025
  \item Para cada trecho, o \textbf{professor} recebe um prompt solicitando raciocínio
        passo a passo (\textit{Chain-of-Thought}) e resposta final em formato JSON
  \item O resultado é salvo em \texttt{distillation\_dataset.jsonl} com os campos
        \texttt{instruction}, \texttt{input}, \texttt{output}, \texttt{reasoning} e \texttt{answer}
\end{enumerate}

\begin{exemplobox}[Exemplo de entrada do dataset sintético]
\textbf{Instrução:} Qual é o CEP da sede da Prefeitura de Aroazes?\\[0.3em]
\textbf{Documento:} ESTADO DO Piauí Prefeitura Municipal DE AROAZES C.N.P.J.:
06.554.984/0001-39 Av. 27 de Fevereiro, 691, Centro CEP: 64310-000\\[0.3em]
\textbf{Reasoning (professor):} To find the CEP of the Aroazes City Hall, I will
identify the relevant field in the document: `CEP: 64310-000'. The CEP is a
Brazilian postal code format (XXXXX-XXX). Therefore, the answer is 64310-000.\\[0.3em]
\textbf{Answer (professor):} 64310-000
\end{exemplobox}

\subsection{Técnica de Destilação: Chain-of-Thought Distillation}

A abordagem utilizada segue o paradigma de \textbf{destilação por Chain-of-Thought
(CoT)}, onde o aluno aprende não apenas as respostas finais do professor, mas também
o raciocínio intermediário. Isso é distinto da destilação clássica por KL-Divergence
de logits (Hinton et al., 2015) e se aproxima da abordagem de \textit{Instruction
Tuning} com dados gerados por LLMs maiores.

\subsection{Hiperparâmetros de Treinamento}

\begin{table}[H]
\centering
\caption{Hiperparâmetros utilizados no fine-tuning do aluno}
\begin{tabular}{ll}
\toprule
\textbf{Parâmetro} & \textbf{Valor} \\
\midrule
Método de fine-tuning   & LoRA (Low-Rank Adaptation) \\
Rank LoRA ($r$)         & 16 \\
Alpha LoRA              & 32 \\
Módulos alvo            & \texttt{q\_proj}, \texttt{v\_proj} \\
Dropout LoRA            & 0.05 \\
Épocas                  & 3 \\
Batch efetivo           & 16 (2 $\times$ 8 grad. accum.) \\
Learning rate           & $2 \times 10^{-4}$ \\
Scheduler               & Linear com warmup \\
Comprimento máximo      & 512 tokens \\
Precisão                & fp16 \\
\bottomrule
\end{tabular}
\end{table}

\subsection{Hardware Utilizado}

\begin{itemize}
  \item \textbf{GPU:} NVIDIA GeForce RTX 4050 (6 GB VRAM)
  \item \textbf{CPU:} Intel Core i5 13ª geração
  \item \textbf{RAM:} 16 GB
  \item \textbf{SO:} Windows 11 + Python 3.12
\end{itemize}

% ═════════════════════════════════════════════════════════════════════════════
\section{Metodologia de Avaliação}
% ═════════════════════════════════════════════════════════════════════════════

\subsection{Avaliação Qualitativa: Benchmark de 100 Perguntas}

O benchmark consiste em 100 pares \texttt{(instrução, contexto, resposta\_esperada)}
extraídos do DOMPI-2025, cobrindo tipos de pergunta como:
extração de CNPJ, endereços, valores de contrato, números de decreto, modalidades de
licitação e obrigações legais.

A avaliação de acerto usou correspondência textual normalizada:
\begin{enumerate}
  \item Normalização: minúsculas, remoção de pontuação e espaços extras
  \item Critério 1 -- \textit{Substring}: gabarito contido na resposta ou vice-versa
  \item Critério 2 -- \textit{Word Overlap}: $\geq 70\%$ das palavras do gabarito
        presentes na resposta
\end{enumerate}

\subsection{Avaliação Complementar: LLM-as-Judge}

Como complemento à métrica automática, as 100 respostas do aluno (antes e depois da
destilação) foram também avaliadas por um \textbf{LLM juiz} (\textit{LLM-as-Judge},
protocolo de \textit{Zheng et al.}, 2023), que recebeu a pergunta, o gabarito e a
resposta do aluno, retornando \textbf{CORRETO}/\textbf{INCORRETO} com justificativa
curta -- tolerando diferenças de formatação e redação, mas não contradições ou dados
inventados (\textit{alucinações}).

Essa abordagem é mais robusta que o \textit{substring matching}: reconhece respostas
corretas com redação diferente do gabarito e, ao mesmo tempo, penaliza respostas que
contêm o trecho certo mas vêm acompanhadas de alucinações -- casos em que a métrica
automática erraria em direções opostas.

% ═════════════════════════════════════════════════════════════════════════════
\section{Resultados}
% ═════════════════════════════════════════════════════════════════════════════

\subsection{Métricas Quantitativas}

\begin{table}[H]
\centering
\caption{Comparação das métricas quantitativas antes e depois da destilação}
\begin{tabular}{lccc}
\toprule
\textbf{Modelo} & \textbf{Perplexidade} & \textbf{Entropia Cruzada} & \textbf{Acurácia Tokens} \\
\midrule
\rowcolor{lightblue}
Professor (3B)    & 4,50 & 1,5050 & 68,08\% \\
Aluno ANTES (1.5B) & 5,49 & 1,7022 & 65,65\% \\
\rowcolor{lightgreen}
Aluno DEPOIS (1.5B) & \textbf{4,23} & \textbf{1,4434} & \textbf{68,65\%} \\
\midrule
\multicolumn{4}{l}{\textit{Variação do aluno (antes $\rightarrow$ depois):}} \\
$\Delta$ & $-1,26\ (-22{,}9\%)$ & $-0,2588\ (-15{,}2\%)$ & $+2,98$ pp \\
\multicolumn{4}{l}{\textit{Gap aluno-depois vs. professor:}} \\
$\Delta$ & $-0,27$ & $-0,0616$ & $+0,57$ pp \\
\bottomrule
\end{tabular}
\end{table}

\begin{successbox}[Resultado-chave das métricas quantitativas]
O aluno destilado \textbf{superou o professor} em todas as três métricas:
PPL $4{,}23 < 4{,}50$, EC $1{,}4434 < 1{,}5050$ e Acurácia $68{,}65\% > 68{,}08\%$.
Isso indica que o fine-tuning especializou o aluno no domínio do dataset (DOMPI-2025),
tornando-o mais eficiente que o próprio professor nesse subconjunto.
\end{successbox}

\subsection{Benchmark Qualitativo: Acertos nas 100 Perguntas (Substring/Word Overlap)}

\begin{table}[H]
\centering
\caption{Resultado do benchmark qualitativo (100 perguntas) -- critério automático}
\begin{tabular}{lcc}
\toprule
\textbf{Métrica} & \textbf{Quantidade} & \textbf{\%} \\
\midrule
Professor acertou              & 55 & 55\% \\
Aluno ANTES acertou            & 36 & 36\% \\
\rowcolor{lightgreen}
Aluno DEPOIS acertou           & \textbf{42} & \textbf{42\%} \\
\midrule
Convergência (ambos acertaram) & 28 & 28\% \\
Aluno acertou, professor errou & 14 & 14\% \\
Professor acertou, aluno errou & 27 & 27\% \\
\midrule
\rowcolor{lightgreen}
Aluno melhorou após destilação & 16 & 16\% \\
\rowcolor{lightred}
Aluno piorou após destilação   & 10 & 10\% \\
\bottomrule
\end{tabular}
\end{table}

O aluno passou de \textbf{36\%} para \textbf{42\%} de acertos (+6 pontos percentuais),
melhorando em 16 questões com apenas 10 regressões. O gap em relação ao professor
(55\% vs 42\%) reflete a diferença de capacidade entre os modelos (3B vs 1.5B).

\subsection{Benchmark Qualitativo: Acertos nas 100 Perguntas (LLM-as-Judge)}

Para verificar a robustez do resultado acima, as mesmas 100 respostas do aluno
(antes e depois) foram reavaliadas por um LLM juiz, conforme descrito na
Seção~3.2. A Tabela~\ref{tab:llm-judge} resume o resultado.

\begin{table}[H]
\centering
\caption{Resultado do benchmark qualitativo (100 perguntas) -- avaliação por LLM-as-Judge}
\label{tab:llm-judge}
\begin{tabular}{lcc}
\toprule
\textbf{Métrica} & \textbf{Quantidade} & \textbf{\%} \\
\midrule
\rowcolor{lightgreen}
Aluno ANTES acertou (juiz)      & \textbf{69} & \textbf{69\%} \\
Aluno DEPOIS acertou (juiz)     & 62 & 62\% \\
\midrule
Convergência -- ambos corretos    & 52 & 52\% \\
Convergência -- ambos incorretos  & 21 & 21\% \\
Aluno melhorou após destilação   & 10 & 10\% \\
\rowcolor{lightred}
Aluno piorou após destilação     & \textbf{17} & \textbf{17\%} \\
\bottomrule
\end{tabular}
\end{table}

\begin{warningbox}[Resultado-chave da avaliação por LLM-as-Judge]
Ao contrário do que sugere a métrica automática de \textit{substring}, a avaliação
semântica por LLM-as-Judge indica que o aluno \textbf{piorou} após a destilação:
de \textbf{69\%} de acertos antes para \textbf{62\%} depois ($-7$ pontos percentuais),
com \textbf{17} questões que regrediram contra apenas \textbf{10} que melhoraram.
Em \textbf{52} das 100 perguntas o julgamento não mudou (ambas as respostas corretas)
e em \textbf{21} ambas seguiram incorretas.
\end{warningbox}

As principais causas das regressões foram: \textbf{alucinações} introduzidas após o
fine-tuning (detalhes inventados, como um endereço ou CNPJ inexistente no
documento-fonte); \textbf{perda de informações} corretas presentes na resposta
``antes'' (o modelo ``depois'' passou a negar tê-las); e \textbf{contradições
internas} (ex.: citar ``R\$~25.000,00'' e, na mesma frase, escrever por extenso
``vinte e quatro mil reais''). Já as 10 melhoras genuínas ocorreram quando a resposta
``depois'' extraiu corretamente um dado ausente ou incompleto na resposta ``antes''
(ex.: matrícula corrigida de ``1234-'' para ``1234-5'').

\subsection{Comparação entre os Dois Métodos de Avaliação}

\begin{table}[H]
\centering
\caption{Comparação: métrica automática (substring/overlap) vs. LLM-as-Judge}
\begin{tabular}{lcc}
\toprule
\textbf{Método de avaliação} & \textbf{Aluno ANTES} & \textbf{Aluno DEPOIS} \\
\midrule
Substring / Word Overlap & 36\% & 42\% ($+6$ pp) \\
LLM-as-Judge             & 69\% & 62\% ($-7$ pp) \\
\bottomrule
\end{tabular}
\end{table}

Os dois métodos divergem na \textbf{direção} do efeito: o \textit{substring}, mais
rígido quanto à forma, subestima o desempenho geral mas indica melhora líquida; o
LLM-as-Judge, ao avaliar o conteúdo semântico, reconhece mais acertos em ambos os
momentos mas aponta piora líquida -- sugerindo que parte do ganho medido pela métrica
automática vem de respostas ``depois'' mais curtas e diretas, que às vezes escondem
alucinações detectáveis apenas por um avaliador semântico.

% ═════════════════════════════════════════════════════════════════════════════
\section{Discussão}
% ═════════════════════════════════════════════════════════════════════════════

\subsection{Houve Transferência de Conhecimento?}

\begin{destaquebox}[Resposta objetiva ao enunciado]
\textbf{Sim, houve transferência de conhecimento}, mas com evidências mistas sobre a
sua qualidade. As métricas quantitativas confirmam a transferência de forma
inequívoca (PPL $-22{,}9\%$, EC $-15{,}2\%$, Acurácia $+2{,}98$ pp -- aluno superou o
professor). Já ao nível de \textbf{corretude semântica} (LLM-as-Judge), o aluno
\textbf{piorou} (69\% $\rightarrow$ 62\%), enquanto pela métrica de
\textit{substring} ele \textbf{melhorou} (36\% $\rightarrow$ 42\%): menor
perplexidade/entropia cruzada não implica, portanto, respostas mais corretas ou
confiáveis.
\end{destaquebox}

\subsection{Por que o Aluno Superou o Professor nas Métricas?}

O resultado é contra-intuitivo, mas explicável: as métricas são calculadas sobre o
\textbf{mesmo dataset} usado no treino do aluno (viés \textit{in-distribution}), e o
fine-tuning com LoRA \textbf{especializou} o aluno no estilo de resposta do professor
nesse domínio específico -- enquanto o professor, genérico, não está especializado
nesse formato. Em um cenário \textit{out-of-distribution} o professor ainda supera o
aluno, como confirma o benchmark qualitativo (55\% vs 42\%, métrica \textit{substring}).

\subsection{Por que a Métrica Automática e o LLM-as-Judge Divergem?}

A divergência (Seção~4.4) ocorre porque o \textit{substring} recompensa respostas
curtas que contenham literalmente o gabarito, sem penalizar alucinações ao lado da
informação correta; já o LLM-as-Judge penaliza tanto a ausência do dado certo quanto
informações fabricadas. Tudo indica que o fine-tuning ensinou o aluno a produzir
respostas num \textbf{formato} mais parecido ao do gabarito (o que ajuda no
\textit{substring}), mas também introduziu maior propensão a alucinar -- o que a
avaliação semântica captura e a automática não. Isso reforça a recomendação de
priorizar avaliações semânticas (LLM-as-Judge ou revisão humana) para conclusões
definitivas sobre corretude.

\subsection{Limitações e Observações}

\begin{itemize}
  \item \textbf{Respostas em inglês do professor:} o Qwen2.5-3B-Instruct respondeu
        em inglês em vários casos, reduzindo artificialmente seus acertos no
        \textit{substring matching}; um prompt mais diretivo mitigaria isso.
  \item \textbf{Tamanho do dataset:} 100 exemplos é pouco para destilação; 1.000+
        (como na Questão 2) dariam resultados mais robustos.
  \item \textbf{Avaliação por substring:} métrica aproximada, como evidenciado pela
        divergência com o LLM-as-Judge (Seções~4.4 e 5.3).
  \item \textbf{Viés do LLM-as-Judge:} o avaliador também é um modelo de linguagem e
        pode ter vieses próprios; o ideal seria complementar com revisão humana.
\end{itemize}

% ═════════════════════════════════════════════════════════════════════════════
\section{Conclusão}
% ═════════════════════════════════════════════════════════════════════════════

Este trabalho demonstrou a viabilidade da destilação de conhecimento por
\textit{Chain-of-Thought} entre modelos da família Qwen2.5, aplicada ao domínio
de documentos oficiais dos municípios do Piauí.

Os principais resultados foram:

\begin{itemize}
  \item \textbf{Transferência confirmada nas métricas quantitativas:} PPL $-22{,}9\%$
        e acurácia de tokens $+2{,}98$ pp após a destilação
  \item \textbf{Especialização no domínio:} o aluno destilado superou o professor
        nas métricas quantitativas
  \item \textbf{Resultado qualitativo misto:} pelo \textit{substring}, o aluno
        melhorou ($+6$ pp); pelo LLM-as-Judge, piorou ($-7$ pp, 17 regressões x 10
        melhoras), em grande parte por alucinações introduzidas no fine-tuning
  \item \textbf{Eficiência:} o aluno (1,5B) alcançou desempenho comparável ao
        professor (3B) treinando via LoRA apenas $\approx$0,5\% dos parâmetros
\end{itemize}

A abordagem confirma que modelos menores podem ser especializados via destilação com
dados sintéticos, reduzindo custos de inferência. Porém, a divergência entre os
métodos de avaliação alerta para o risco de conclusões precipitadas baseadas apenas
em correspondência textual: ganhos de formato podem mascarar perda de confiabilidade
factual, que só uma avaliação semântica (ou humana) captura.

% ═════════════════════════════════════════════════════════════════════════════
\section*{Referências}
% ═════════════════════════════════════════════════════════════════════════════

\begin{enumerate}[label={[\arabic*]}]
  \item \label{dompi2025} Portelaa, G. et al. (2025). \textit{DOMPI-2025: Dataset dos
        Diários Oficiais dos Municípios do Piauí}.
        \url{https://huggingface.co/datasets/gutoportelaa/DOMPI-2025}
  \item \label{llmjudge} Zheng, L. et al. (2023). \textit{Judging LLM-as-a-Judge with
        MT-Bench and Chatbot Arena}. arXiv:2306.05685.
\end{enumerate}

\end{document}